%% Master CV - Priyadarshi Roy
%% Comprehensive overview of all competencies, experience, and achievements.
%% Use as a master reference when tailoring role-specific CVs.
%%
%% Compile with: lualatex main_example.tex
%% (pdflatex often fails on modern MiKTeX with fontawesome5 font-expansion errors.)

\documentclass[11pt,a4paper,sans]{moderncv}
\moderncvstyle{banking}
\moderncvcolor{blue}

% Force the name and section headings to render in moderncv blue (color1).
% Default banking leaves them black: moderncvstylebanking.sty's \colorlet
% copies (not aliases) the pre-scheme accent colour, so the name colours are
% frozen before \moderncvcolor runs. Re-let them after. \namefont is the hook
% every name-style macro routes through, so this also works on moderncv 2.3.1
% (Debian/Ubuntu apt), which has no \firstnamestyle/\lastnamestyle at all.
\renewcommand*{\namefont}{\fontsize{34}{36}\bfseries\upshape}
\colorlet{firstnamecolor}{color1}
\colorlet{lastnamecolor}{color1}
\colorlet{namecolor}{color1}
\renewcommand*{\sectionstyle}[1]{{\sectionfont\color{color1}#1}}

\usepackage[utf8]{inputenc}
% pdflatex fallback only (the documented engine is lualatex, which skips this
% branch). Without T1 font encoding pdflatex builds accented letters with
% \accent, and the PDF text layer stores them decomposed - `e` + U+0300 rather
% than U+00E8 - so an ATS keyword match on "Genève" fails while the page looks
% right. moderncv 2.5 loads T1 itself under pdflatex; 2.3.1 (Debian/Ubuntu apt)
% does not. \ifpdftex comes from iftex, which every moderncv version loads.
\ifpdftex\usepackage[T1]{fontenc}\fi
% moderncv loads hyperref itself in an \AtEndPreamble hook, so \hypersetup
% must go in an \AtEndPreamble of our own: on moderncv < 2.4 a top-level
% \usepackage{hyperref} clashes with the class's own
% \RequirePackage[unicode]{hyperref}. From 2.4.0 the class passes its options
% through \PassOptionsToPackage instead, which is what removes that clash.
\AtEndPreamble{\hypersetup{
    colorlinks=true,
    linkcolor=blue,
    filecolor=magenta,
    urlcolor=blue,
    pdftitle={Priyadarshi Roy - CV},
    % Keep pdfpagemode=UseNone: this block runs after moderncv's own
    % \AtEndPreamble (moderncv.cls sets pdfpagemode there), so a FullScreen
    % value here would win and open every CV in fullscreen presentation mode.
    pdfpagemode=UseNone,
}}
\usepackage[scale=0.80]{geometry}
\usepackage{import}
\usepackage{needspace}

% personal data
\name{Priyadarshi}{Roy}
\address{Cölbe, Germany}{}{}
\phone[mobile]{+49 1521 7847618}
\email{proy9714@gmail.com}
\extrainfo{\href{https://www.linkedin.com/in/priyadarshi-roy/}{linkedin.com/in/priyadarshi-roy}, \href{https://github.com/proy9714}{github.com/proy9714}}

\begin{document}

\makecvtitle

% ============================================================
%     PROFILE STATEMENT
% ============================================================

\vspace{6pt}
\small{Machine Learning Engineer with 3 years of industry experience building production NLP and Generative AI systems, including enterprise Retrieval-Augmented Generation (RAG) applications, LLM orchestration and evaluation pipelines, and Python REST APIs serving models in production. MSc Data Science student at Philipps-Universität Marburg (thesis in multimodal computer vision, submission November 2026), available immediately.}

% ============================================================
%     CORE COMPETENCIES
% ============================================================

\section{Core Competencies}
\vspace{1pt}
\begin{itemize}

\item \textbf{LLMs and Generative AI}: Retrieval-Augmented Generation (RAG), LangChain, prompt engineering, semantic and hybrid search, Text-to-SQL, agentic workflows, LLM evaluation with RAGAS.

\item \textbf{Machine Learning and NLP}: Python, PyTorch, Hugging Face Transformers, BERT, TensorFlow, Keras, Scikit-Learn; document summarisation, text classification, imbalanced learning.

\item \textbf{Computer Vision and Multimodal AI}: Vision Transformers, vision-language models, human pose estimation, object detection, image retrieval, OpenCV, ASR/TTS.

\item \textbf{Backend and MLOps}: FastAPI, Flask, REST APIs, Docker, Docker Compose, Git, Linux, Ollama.

\item \textbf{Data Engineering}: PostgreSQL, pgvector, MongoDB, MySQL, Apache Spark, Pandas, SQL, Java, C++, JavaScript.

\end{itemize}

% ============================================================
%     PROFESSIONAL EXPERIENCE
% ============================================================

\section{Professional Experience}
\vspace{3pt}
\begin{itemize}

\item{\cventry{Aug 2021 - Aug 2024}{Machine Learning Engineer}{Sprinriver Technology Private Limited}{Kolkata, India}{}{\vspace{1pt}
\begin{itemize}
    \item Developed production NLP systems for automated document understanding, including multi-document summarisation and highlight extraction pipelines that replaced manual review of large document collections.
    \item Built enterprise Retrieval-Augmented Generation (RAG) applications using GPT-3.5 Turbo, embedding-based retrieval, LangChain and vector search, enabling employees to query organisational knowledge in natural language.
    \item Designed modular LLM orchestration workflows and evaluation pipelines, improving the reliability, maintainability and scalability of Generative AI applications.
    \item Developed and deployed Python REST APIs with Flask to serve machine learning models in production, supporting AI-powered features across business applications.
\end{itemize}}}

\end{itemize}

% ============================================================
%     RESEARCH EXPERIENCE
% ============================================================

\section{Research Experience}
\vspace{3pt}
\begin{itemize}

\item{\cventry{2026-present}{Master's Thesis: Computer Vision for Historical Art Analysis}{Philipps-Universität Marburg}{Germany}{}{\vspace{1pt}
\begin{itemize}
    \item Developing multimodal pipelines that combine vision foundation models, human pose estimation, object detection, image embeddings and geometric verification to analyse human figures and interactions in historical paintings.
    \item Investigating domain shift from natural images to artistic media, including scene graph generation and knowledge-based retrieval.
\end{itemize}}}

\vspace{3pt}

\needspace{5\baselineskip}
\item{\cventry{2026}{Independent Scientific Project: ``Noli Me Tangere''}{Philipps-Universität Marburg}{Germany}{}{\vspace{1pt}
\begin{itemize}
    \item Designed a pose-grounded human-human interaction scene-graph pipeline (Grounding DINO detection, pose embeddings, vision-language-model reasoning) and assembled a corpus of about 1,000 artworks from ART Research Network, JSTOR and Prometheus.
\end{itemize}}}

\end{itemize}

% ============================================================
%     PROJECTS
% ============================================================

\needspace{8\baselineskip}
\section{Projects}
\vspace{3pt}
\begin{itemize}

\needspace{5\baselineskip}
\item{\cventry{2025}{Python, FastAPI, LangChain, PostgreSQL, pgvector, Vue.js}{Study Advisor (RAG-Based Chatbot)}{}{}{\vspace{1pt}
\begin{itemize}
    \item Architected a multi-stage RAG pipeline for academic advising with intent classification, query decomposition, Text-to-SQL, semantic retrieval, context summarisation and LLM response generation.
    \item Implemented hybrid retrieval inside PostgreSQL using pgvector (semantic search) and pg\_trgm (fuzzy matching), removing the need for a separate vector database.
    \item Benchmarked Llama 3.1-70B against Gemini 2.5-Flash with RAGAS over 75 questions; Gemini led on faithfulness (0.65 vs.\ 0.44) and answer correctness (0.57 vs.\ 0.47).
\end{itemize}}}

\vspace{3pt}

\needspace{5\baselineskip}
\item{\cventry{2026}{Python, Faster-Whisper, Ollama, Docker Compose}{NLS Voice Assistant}{}{}{\vspace{1pt}
\begin{itemize}
    \item Built a voice assistant (speech recognition, local LLM orchestration, text-to-speech) for weather and multi-turn calendar tasks, with a scenario-based evaluation framework reaching 100\% intent accuracy.
\end{itemize}}}

\vspace{3pt}

\needspace{5\baselineskip}
\item{\cventry{2025}{Python, PyTorch, Hugging Face Transformers, Scikit-Learn}{Multi-Label Emotion Classification}{}{}{\vspace{1pt}
\begin{itemize}
    \item Applied MLSMOTE in BERT-embedding space to address severe label imbalance (2,768 to 4,724 training samples across 28 label combinations), reaching a weighted F1 of 0.81.
\end{itemize}}}

\vspace{3pt}

\needspace{5\baselineskip}
\item{\cventry{2024}{\href{https://github.com/assafelovic/gpt-researcher}{github.com/assafelovic/gpt-researcher}}{Open Source: GPT Researcher (Contributor)}{}{}{\vspace{1pt}
\begin{itemize}
    \item Contributed features, bug fixes and code improvements to an open-source autonomous research agent (LLM orchestration, RAG, web retrieval) through reviewed pull requests.
\end{itemize}}}

\vspace{3pt}

\item \textbf{Earlier projects}: Distributed inclusion-dependency discovery (2024; Java, Akka, Spark); NLP models for fake news classification (2021; TensorFlow, Keras); CNN hand-gesture recognition (2019; OpenCV).

\end{itemize}

% ============================================================
%     EDUCATION
% ============================================================

\section{Education}
\vspace{1pt}
\begin{itemize}

\needspace{4\baselineskip}
\item{\cventry{2024-present}{MSc Data Science}{Philipps-Universität Marburg}{Germany}{}{\vspace{1pt}
In progress, thesis submission November 2026. Computer Vision, Natural Language Systems, ML with Graphs.
}}

\vspace{3pt}

\item{\cventry{2019-2021}{MSc Computer Science}{St. Xavier's College (Autonomous)}{Kolkata, India}{}{\vspace{1pt}
CGPA 8.17/10, First Class.
}}

\vspace{3pt}

\item{\cventry{2016-2019}{BSc Computer Science}{St. Xavier's College (Autonomous)}{Kolkata, India}{}{\vspace{1pt}
CGPA 7.99/10, First Class.
}}

\end{itemize}

% ============================================================
%     CERTIFICATIONS
% ============================================================

\section{Certifications}
\vspace{1pt}
\begin{itemize}
\item {LangChain for LLM Application Development; LangChain: Chat with Your Data (DeepLearning.AI); NLP with Classification and Vector Spaces; Build Basic GANs (Coursera); Python and Computer Vision Bootcamp (ShapeAI).}
\end{itemize}

% ============================================================
%     LANGUAGES
% ============================================================

\section{Languages}
\vspace{1pt}
\begin{itemize}
\item {English (C2), Bengali (native), Hindi (C2), Japanese (JLPT N3), German (A2).}
\end{itemize}

% ============================================================
%     REFERENCES
% ============================================================

\section{References}
\vspace{1pt}
\begin{itemize}
\item{Available upon request.}
\end{itemize}

\end{document}
